% ATTEMPT_STATUS: unresolved
% TOP_PROBLEM: {"problemId": "problem.riemann-hypothesis", "canonicalTitle": "Riemann Hypothesis", "releaseRank": 2, "primaryDomain": "number_theory_arithmetic_geometry", "primaryDomainLabel": "Number theory & arithmetic geometry", "displayStatus": "open", "releaseStatus": "open"}
% !TeX program = lualatex
% Complete problem section copied verbatim from research_batch01.tex.
\documentclass[11pt]{article}
\usepackage[margin=25mm]{geometry}
\usepackage{fontspec}
\setmainfont{Latin Modern Roman}
\usepackage{amsmath,amssymb,mathtools}
\usepackage{unicode-math}
\setmathfont{Latin Modern Math}
\usepackage{enumitem,needspace}
\usepackage{xurl,hyperref}
\hypersetup{colorlinks=true,urlcolor=blue}
\setcounter{secnumdepth}{1}
\setlength{\parindent}{0pt}
\setlength{\parskip}{5pt}
\setlength{\emergencystretch}{3em}
\setlist[itemize]{leftmargin=3em,itemsep=5pt,topsep=4pt,parsep=0pt}
\allowdisplaybreaks
\newcommand{\N}{\mathbb{N}}
\newcommand{\Z}{\mathbb{Z}}
\newcommand{\Q}{\mathbb{Q}}
\newcommand{\R}{\mathbb{R}}
\newcommand{\C}{\mathbb{C}}
\newcommand{\F}{\mathbb{F}}
\newcommand{\A}{\mathbb{A}}
\newcommand{\PP}{\mathbb P}
\newcommand{\E}{\mathbb{E}}
\newcommand{\eps}{\varepsilon}
\newcommand{\dd}{\,\mathrm d}
\newcommand{\sref}[2]{\hyperlink{source:#1:#2}{[S#2]}}
\newcommand{\eref}[2]{\hyperlink{extra:#1:#2}{[E#2]}}
% Turn the next switch off for a shorter reading copy. The quotations preserve
% every exactTarget field, including qualifications omitted from a short summary.
\newif\ifshowcatalogue
\showcataloguetrue
\newcommand{\cataloguescope}[1]{\ifshowcatalogue
 \paragraph{Exact scope in the supplied catalogue.}
 \begin{quote}\small #1\end{quote}\fi}
\begin{document}
\setcounter{section}{1}
% ================================================================
% problemId: problem.riemann-hypothesis
% source releaseRank: 2; source releaseStatus: open
% exactTarget SHA-256: 43a61dff12aa4263f3e68f91192ddf29e0218c93b508ad4fad99b54d1df28674
\clearpage
\section{\texorpdfstring{Riemann Hypothesis}{Riemann Hypothesis}}\label{2}
\subsection{Definitions and mathematical statement}
For $s\in\C$ with $\Re s>1$, define
\[
 \zeta(s)=\sum_{n=1}^{\infty}n^{-s}=\prod_{p\text{ prime}}(1-p^{-s})^{-1}.
\]
Use its meromorphic continuation to $\C$, whose only pole is at $s=1$. The assertion is
\[
 \forall\rho\in\C,\qquad \bigl(0<\Re\rho<1\ \text{and}\ \zeta(\rho)=0\bigr)
 \ \Longrightarrow\ \Re\rho=\tfrac12.
\]
The open critical strip excludes the negative even integer zeros. The assertion concerns locations, not multiplicities, of zeros.
\par\smallskip\textit{Formulation and concept references:} \sref{2}{1}.
\cataloguescope{For the meromorphic continuation of the Riemann zeta function zeta(s)=sum\_\{n>=1\} n\textasciicircum{}(-s), initially defined for Re(s)>1, every zero s in the open critical strip 0<Re(s)<1 has Re(s)=1/2. This is the Riemann Hypothesis for zeta alone; it does not assert simplicity of its zeros.}
\subsection{Short English statement}
Are all zeta zeros between the vertical lines zero and one located on the halfway line?
% BEGIN RESEARCH ATTEMPT -- batch 1, problem 2, 24 September 2026
\subsection{Research attempt}

\paragraph{Current frontier and status distinctions.}
The exact target is the zero-location assertion for $\zeta$ alone, with
multiplicities unrestricted \sref{2}{1}. An established computational theorem
of Platt--Trudgian verifies it throughout $0<\Im\rho\le 3\cdot10^{12}$,
using interval arithmetic and a count of all zeros, not merely a list of zeros
found on the line \eref{2}{1}. On the analytic side, Guth--Maynard's work,
initiated in 2024 and published in 2026, gives the zero-density estimate
\[
 N(\sigma,T)\le T^{15(1-\sigma)/(3+5\sigma)+o(1)},
 \qquad N(\sigma,T)=\#\{\rho:\Re\rho\ge\sigma,\ |\Im\rho|\le T\},
\]
for fixed $1/2<\sigma<1$ as $T\to\infty$, with multiplicities counted
\eref{2}{2}, Theorem~1.2. Such an upper bound does
not assert that the exceptional set is empty.

A substantial recent development in the supplied bibliography must not be
omitted: Lamzouri's September 2026 manuscript states the unconditional bound
\[
 \liminf_{T\to\infty}\frac{N_0^{\rm s}(T)}{N(T)}
 \ge \frac32-\frac1{\sqrt2}\cot\!\left(\frac1{\sqrt2}\right)
 =0.6725007\ldots,
\]
where $N_0^{\rm s}$ counts simple zeros on the line and $N$ counts all upper-half
strip zeros with multiplicity \sref{2}{3}, \eref{2}{9}, Theorem~1.1.
The associated September 3 lecture is also listed by the host institution
\sref{2}{4}. The theorem statement and relevant mechanism were inspected;
the complete analytic proof and its formal-certificate dependencies were not
independently verified in this run. It is recorded as a recent research result
at manuscript stage, not as an established resolution of RH. Even an
asymptotic proportion of one would not, by itself, exclude finitely many
exceptions.

Two established results sharpen the positivity route below. Griffin--Ono--
Rolen--Zagier prove eventual hyperbolicity of the central Jensen polynomials
for every fixed degree; the threshold can depend on that degree
\eref{2}{3}, Theorem~1. A recent, unaudited August 2026 manuscript of Holland
states hyperbolicity in the joint range $n^3\log^2(n+2)\ge Kd^5$ for an absolute
$K>0$, still not for every pair $(d,n)$ \eref{2}{6}, Theorem~1.1.
Rodgers--Tao prove that the de Bruijn--Newman constant is nonnegative; since
RH is equivalent to its nonpositivity, the required endpoint is exactly zero
\eref{2}{4}. These are selected frontier results relevant to this attempt,
not a claim that all current research is exhausted by them.

\paragraph{Attack map.}
A heat-flow argument would have to establish real-rootedness down to the exact
endpoint, not merely at positive deformation parameters \eref{2}{4}.
The new pair-correlation approach would need additional control that detects
every possible exception, beyond a global proportion estimate
\eref{2}{9}, Sections~2--3. A factorization argument would need to show why the
functional equation fixes individual zero factors instead of permuting them;
the first stress test below shows that symmetry alone cannot do this.
The selected route combines local Taylor data, logarithmic derivatives,
and polynomial filters that isolate hypothetical nonreal inverse zeros.
Its precise missing lemma is positivity of one explicit quadratic form on
\emph{all} real polynomials. Unlike a numerical search over zeros, the form
can be defined without assuming where any zeros lie.

The calculations and auxiliary proofs below were developed in this notebook.
Positivity criteria and moment methods are classical; no historical novelty
is asserted for the resulting equivalence. The aim is to expose a usable
intermediate target and determine whether the available positivity actually
establishes it.

\paragraph{Attempt A: source-audit obstruction, not a zeta counterexample.}
Theorem~1 of the existing proof claim \sref{2}{5}, version~19,
Section~3, passes from factor matching in its equation~(5) to same-index
matching in equation~(7), invoking uniqueness of zero multiplicities. That
inference is invalid: multiplicity is attached to a zero, but equal
multiplicities do not distinguish different zeros. This can be tested within
the infinite-product class used there, without appealing to another claimed
solution. Fix $a=1/4$ and put
\begin{equation}\label{ra:2:symmetry-test}
 A(s)=\prod_{n=1}^{\infty}
 \left(1+\frac{(s-a)^2}{n^2}\right)
 \left(1+\frac{(s-(1-a))^2}{n^2}\right).
\end{equation}
On each compact set, the sum of the absolute deviations of these factors from
one is bounded by a constant times $\sum n^{-2}$, so the product converges
locally uniformly to a nonzero entire function. Away from the displayed
factor zeros its product is nonzero; hence its zeros are exactly
$a\pm in$ and $1-a\pm in$, all simple. Their real parts lie strictly between
zero and one, and the sum of inverse squared imaginary parts is finite.
Reflection $s\mapsto1-s$ exchanges the two factors at each $n$, so
$A(s)=A(1-s)$, although none of its zeros lies on the critical line.
Each quadratic factor is irreducible over $\R$. Thus even within this
convergent, strip-confined, multiplicity-preserving setting, reflection may
exchange two distinct conjugate-pair factors of the same multiplicity.

This explicitly defeats the same-index inference in the cited argument.
It does \emph{not} disprove RH: $A$ is not the completed zeta function, and no
Euler-product identity for it is asserted. In particular, merely using an
Euler product to confine zeros to a strip does not repair that inference;
more arithmetic information would have to enter the proof.

\paragraph{Attempt B: define the moment form without locating zeros.}
To avoid a clash with the $t$-variable convention in \sref{2}{1}, throughout
this attempt write
\[
 \xi(s)=\tfrac12s(s-1)\pi^{-s/2}\Gamma(s/2)\zeta(s),\qquad
 X(z)=\frac{\xi(1/2+z)}{\xi(1/2)},\qquad
 F(w)=\sum_{n=0}^{\infty}a_nw^n,\quad X(z)=F(z^2).
\]
The functional equation makes $X$ even, so this definition of $F$ is entire
and has no square-root branch choice. Also $a_0=1$. There is no real strip zero:
for $0<s<1$ the alternating series
$\eta(s)=\sum_{n\ge1}(-1)^{n-1}n^{-s}$ is strictly positive (group adjacent
positive differences), whereas $1-2^{1-s}<0$ and
$\zeta(s)=\eta(s)/(1-2^{1-s})$. In particular $\xi(1/2)>0$.
The completed-function properties and zero correspondence are as in
\sref{2}{1}; the growth bound
$|\xi(s)|\le\exp(O(|s|\log(2+|s|)))$ is recorded in
\eref{2}{7}, Lemma~1.

Define real numbers $m_k$ solely by the convergent expansion near zero
\begin{equation}\label{ra:2:moment-definition}
 \frac{F'(w)}{F(w)}=\sum_{k=0}^{\infty}(-1)^km_kw^k.
\end{equation}
Equating coefficients in $F'=F(F'/F)$ gives the exact recurrence
\begin{equation}\label{ra:2:recurrence}
 m_k=(-1)^k\left((k+1)a_{k+1}
       -\sum_{j=1}^{k}a_j(-1)^{k-j}m_{k-j}\right).
\end{equation}
For example,
\[
 m_0=a_1,\qquad m_1=a_1^2-2a_2,\qquad
 m_2=a_1^3-3a_1a_2+3a_3.
\]
Set $\mathcal L(t^k)=m_k$, extending linearly to $\R[t]$, and let
\begin{equation}\label{ra:2:hankel}
 H_D=(m_{i+j})_{0\le i,j\le D}.
\end{equation}
Thus $c^{\mathsf T}H_Dc=\mathcal L(p^2)$ for
$p(t)=\sum_{i=0}^D c_it^i$. The unresolved estimate is
\begin{equation}\label{ra:2:missing-positivity}
 \mathcal L(p^2)\ge0\qquad\text{for every }p\in\R[t].
\end{equation}
We next prove carefully that this estimate would address the full target,
not merely a density-one or bounded-height version.

\medskip
\textbf{Unconditional zero expansion.}
Let the distinct zeros of $F$ be $w_j$, with multiplicities $d_j$, and put
$\lambda_j=-1/w_j$. Equivalently, $w_j=(\rho-1/2)^2$, with each pair
$\{\rho,1-\rho\}$ represented once. The standard zero count
$O(T\log T)$ implies $\sum_jd_j|\lambda_j|<\infty$; finitely many bounded
zeros cause no problem because $F(0)=1$. Hadamard factorization applied to
$X$ pairs its zeros $b,-b$ into
$(1-z/b)e^{z/b}(1+z/b)e^{-z/b}=1-z^2/b^2$.
Absolute convergence of the canonical product permits this pairing.
The remaining exponential is $e^{Bz}$, and evenness forces $B=0$.
Consequently
\begin{equation}\label{ra:2:product}
 F(w)=\prod_j(1+\lambda_jw)^{d_j},\qquad
 m_k=\sum_jd_j\lambda_j^{k+1}.
\end{equation}
The factorization theorem and growth/zero-count prerequisites are in
\eref{2}{7}, Section~1; see also \eref{2}{8}, Section~3.
For $|w|\sup_j|\lambda_j|<1/2$, both the logarithmic-derivative sum and
its geometric expansions converge absolutely, bounded by a constant times
$\sum_jd_j|\lambda_j|$. This justifies the coefficient interchange in
\eqref{ra:2:product}. It does not assume RH.

The $\lambda_j$ are closed under conjugation, with equal multiplicities.
They have no accumulation point except zero. Since a nontrivial zeta zero
has nonzero imaginary part, elementary squaring gives
\begin{equation}\label{ra:2:lambda-location}
 \lambda_j\in\R\quad\Longleftrightarrow\quad
 \Re\rho=\tfrac12,
 \qquad\text{and then }\lambda_j=(\Im\rho)^{-2}>0.
\end{equation}
For every fixed polynomial, absolute convergence yields
\begin{equation}\label{ra:2:Lzero}
 \mathcal L(p^2)=\sum_j d_j\lambda_jp(\lambda_j)^2.
\end{equation}
The square here is \emph{not} an absolute square unless its argument is real.
Mistaking these expressions would assume precisely the desired conclusion.

\medskip
\textbf{Lemma 2.1 (polynomial isolation of a hypothetical off-axis pair).}
If some $\lambda_*$ in \eqref{ra:2:product} is nonreal, there is a real
polynomial $q$ for which $\mathcal L(q^2)<0$.

\textit{Proof.}
Put $r=|\lambda_*|>0$, and let $S$ consist of all distinct $\lambda_j$
with $|\lambda_j|\ge r$, except $\lambda_*$ and $\overline{\lambda_*}$.
This is a finite conjugation-invariant set. Therefore
\[
 P(t)=\prod_{\lambda\in S}(t-\lambda)\in\R[t],
 \qquad P(\lambda_*)\ne0.
\]
An empty product is one. Choose a complex number $\eta$ with
$\lambda_*\eta^2=-1$. For each integer $k\ge0$ choose real $u_k,v_k$ solving
\[
 u_k+v_k\lambda_*
       =\frac{\eta}{\lambda_*^kP(\lambda_*)},
 \qquad q_k(t)=t^kP(t)(u_k+v_kt).
\]
The real-linear map $(u,v)\mapsto u+v\lambda_*$ is invertible because
$\Im\lambda_*\ne0$. Its inverse has fixed finite norm, so
$|u_k|+|v_k|\le C_*r^{-k}$, with $C_*$ independent of $k$.
We have $q_k(\lambda_*)=\eta$ and
$q_k(\overline{\lambda_*})=\overline\eta$. The contribution of this pair
to \eqref{ra:2:Lzero} is exactly $-2d_*$, where $d_*$ is its common
multiplicity. All terms in $S$ vanish.

The remaining inverse zeros have modulus at most some $R<r$: otherwise
infinitely many would accumulate at modulus $r>0$. If there are no remaining
zeros the proof is already complete. On $|t|\le R$, the fixed polynomial
$P$ is bounded, hence
\[
 |q_k(t)|\le C'_* (R/r)^k.
\]
It follows that the absolute value of the remaining contribution is at most
\begin{equation}\label{ra:2:filter-tail}
 (C'_*)^2(R/r)^{2k}\sum_{\lambda_j\notin S\cup\{\lambda_*,\overline{\lambda_*}\}}
                  d_j|\lambda_j|\ \longrightarrow\ 0.
\end{equation}
Thus $\mathcal L(q_k^2)=-2d_*+o(1)<0$ for all sufficiently large $k$.
All operations used a finite polynomial and an absolutely summable tail.
\hfill$\square$

\medskip
\textbf{Consequence (an exact positivity reduction).}
RH is equivalent to $H_D\succeq0$ for every $D\ge0$. If RH holds, then
\eqref{ra:2:lambda-location} and \eqref{ra:2:Lzero} give positivity
term by term. In fact every $H_D$ is positive definite, since there are
infinitely many distinct positive $\lambda_j$ and a nonzero polynomial has
only finitely many roots. Conversely, if RH fails, Lemma~2.1 supplies a
finite-degree polynomial violating positivity of one $H_D$. This argument
does not impose simplicity of any zero. Its usefulness remains conditional:
the proof constructs the negative polynomial \emph{assuming a nonreal
inverse zero has already been specified}; it neither produces such a zero
for $\zeta$ nor proves positivity from arithmetic.

\paragraph{Attempt C: can the positive theta kernel prove the missing estimate?}
In the normalization of \eref{2}{4}, equations~(1)--(3), set
\[
 \Phi(u)=\sum_{n\ge1}
 (2\pi^2n^4e^{9u}-3\pi n^2e^{5u})e^{-\pi n^2e^{4u}},\qquad u\ge0.
\]
Each summand is positive because its prefactor is
$\pi n^2e^{5u}(2\pi n^2e^{4u}-3)>0$. The cited Fourier representation gives
\begin{equation}\label{ra:2:theta}
 X(z)=\frac{\int_0^\infty\Phi(u)\cosh(2zu)\,du}
            {\int_0^\infty\Phi(u)\,du},\qquad
 a_n=\frac{2^{2n}}{(2n)!}
      \frac{\int_0^\infty u^{2n}\Phi(u)\,du}
           {\int_0^\infty\Phi(u)\,du}.
\end{equation}
The super-exponential decay of $\Phi$ dominates
$e^{2|z|u}$ on every compact $z$-set, so differentiation and expansion under
the integral are justified. In particular all $a_n$ are positive.
The attempted bridge is to combine ordinary moment positivity of the measure
$\Phi(u)du$ with recurrence~\eqref{ra:2:recurrence} to prove positivity of
$\mathcal L$. But the factorial reweighting and the logarithmic derivative
are nonlinear operations; ordinary moment positivity does not automatically
survive them.

Here is an exact falsification of that bridge in its proposed generality.
With the rescaled integration variable $v=2u$, take the positive probability
measure with masses $3/4$ at $v=0$ and $1/4$ at $v=1$. Its normalized transform
$\int\cosh(v\sqrt w)\,d\mu(v)$ is
\[
 F_{\rm toy}(w)=\tfrac34+\tfrac14\cosh\sqrt w
     =1+\tfrac18w+\tfrac1{96}w^2+\tfrac1{2880}w^3+\cdots.
\]
Its coefficients are nonnegative and its square-root notation again denotes
an entire power series. The same recurrence gives
\[
 m_0=\tfrac18,\qquad m_1=-\tfrac1{192},\qquad
 m_2=-\tfrac7{7680},\qquad
 \det H_1=-\tfrac{13}{92160}<0.
\]
Indeed, the equation $\cosh z=-3$ has the nonimaginary solutions
$z=\pm\operatorname{arcosh}(3)+(2\ell+1)\pi i$.
This is not a counterexample to RH, nor a replacement of the actual theta
kernel. It proves that positivity of a cosh-integrating measure alone cannot
justify the proposed step. Even smooth nonnegative kernels cannot rescue
that general statement: replace the two atoms by narrow nonnegative smooth
bumps of total masses $3/4$ and $1/4$. Their first six moments converge to
those of the atoms, so the displayed strictly negative determinant persists.
Any successful use of \eqref{ra:2:theta} must therefore exploit an additional,
specific property of this theta kernel.

\paragraph{Attempt D: test whether finite positivity plus zero symmetry could suffice.}
Consider the explicitly defined model
\[
 F_*(w)=\prod_{j\ge1}(1+w/j^2).
\]
It is entire, and its inverse zeros are the distinct positive numbers $j^{-2}$.
Its moment matrices are all positive definite by the same finite-polynomial
argument used above. The compact-uniform convergence also shows that its
Taylor coefficients are positive. Fix a degree $D$ and let $\delta_D>0$
be the smallest eigenvalue of its $H_D$. For $0<a<1/2$, choose $T>a$ and set
\begin{equation}\label{ra:2:perturb}
 \lambda=-\frac1{(a+iT)^2},\qquad
 \widetilde F(w)=F_*(w)(1+\lambda w)(1+\overline\lambda w).
\end{equation}
Then $\Re\lambda>0$ and
$(1+\lambda w)(1+\overline\lambda w)
=1+2\Re\lambda\,w+|\lambda|^2w^2$ has positive coefficients.
The additional zeros of $\widetilde F(z^2)$ are exactly
$z=\pm a\pm iT$. Thus the translated $s$-function remains entire, real under
conjugation, symmetric under $s\mapsto1-s$, and has all zeros in the open
critical strip, but it has an off-line quartet.

Its moment perturbation is exactly
\[
 \Delta m_k=2\Re(\lambda^{k+1}).
\]
For $v=(1,\lambda,\ldots,\lambda^D)^{\mathsf T}$, this gives
\[
 \Delta H_D=\lambda vv^{\mathsf T}
           +\overline\lambda\,\overline v\,\overline v^{\mathsf T},
 \qquad
 \|\Delta H_D\|_2\le
      2|\lambda|\sum_{i=0}^D|\lambda|^{2i}
      \le\frac{2|\lambda|}{1-|\lambda|^2}\quad(|\lambda|<1).
\]
Since $|\lambda|=(T^2+a^2)^{-1}\to0$, choose $T$ so that this last bound
is less than $\delta_D$. The perturbed $H_D$ and all its leading submatrices
remain positive definite despite the off-line quartet. Moreover the
perturbation of any prescribed finite list of moments can be made arbitrarily
small. This is an auxiliary construction, not a modified $\zeta$ Euler product.
It shows exactly why a finite group of successful tests plus these generic
symmetries is inadequate; it does not rule out a finite argument using a
stronger arithmetic invariant.

The polynomial-isolation proof explains the complementary phenomenon:
it eventually detects the quartet, but the degree and coefficient sizes of
its filter depend on the unobserved inverse zeros and on the separation
$R<r$ in \eqref{ra:2:filter-tail}. No uniform small-degree bound was obtained.
Farmer's analysis of information lost under repeated differentiation is a
related warning about extrapolating from Jensen-polynomial asymptotics
\eref{2}{5}, Sections~2--4; it is not a theorem prohibiting all positivity
methods.

\paragraph{Computational checks and reproducibility.}
Exact rational checks in the accompanying script verify the toy fractions.
They also confirm that, for the
finite base $\prod_{j=1}^5(1+w/j^2)$ and the perturbation
$a=1/4$, $T=10^6$, every leading determinant through the $4\times4$ moment
matrix is positive both before and after adding the off-axis quartet.
The corresponding zeros are specified algebraically, not inferred numerically.
In the same finite model, take
$P(t)=\prod_{j=1}^5(t-j^{-2})$ and solve
$u+v\lambda=(a+iT)/P(\lambda)$ for rational $u,v$.
Then $q(t)=P(t)(u+vt)$ annihilates every base inverse zero and satisfies
$\lambda q(\lambda)^2=-1$. Thus
$\mathcal L_{\rm model}(q^2)=-2$, as the code independently verifies using
the rational moment matrix and polynomial coefficients. The polynomial has
degree six: the same model passes the displayed low-degree tests and fails
at a higher degree.

For the actual $\xi$, the script computes its Taylor coefficients at $1/2$
through order $14$ and applies \eqref{ra:2:recurrence}. It does \emph{not}
compute these moments from a list of presumed on-line zeros, which would
silently omit any exceptions. Runs at $70$ and $90$ decimal working digits
agreed to the displayed precision. Illustrative values are
\[
 \begin{aligned}
 m_0&\approx0.02310499311541897079,\\
 m_1&\approx3.717259928526968616\cdot10^{-5},\\
 m_2&\approx1.441739314009732797\cdot10^{-7}.
 \end{aligned}
\]
For numerical conditioning, also form the diagonal congruence
$\widehat H_D=(m_{i+j}/\sqrt{m_{2i}m_{2j}})_{i,j=0}^D$; all required
$m_{2i}$ were positive in these computations. The observed results are
\[
\begin{array}{c|c|c}
 D&\det H_D&\lambda_{\min}(\widehat H_D)\\ \hline
 0&2.3104993115\cdot10^{-2}&1\\
 1&1.9493355548\cdot10^{-9}&0.3559395401\\
 2&2.1728105105\cdot10^{-22}&0.01549146869\\
 3&1.1775817160\cdot10^{-41}&0.0002184026834
\end{array}
\]
These are ordinary high-precision floating-point calculations, not interval
certificates. Repetition at higher precision is a diagnostic, not a proof of
any determinant sign. Even certified signs for this finite table would not
establish the universal condition \eqref{ra:2:missing-positivity}.

\paragraph{Stress test and exact remaining gap.}
The strongest step supplied here is the full finite-witness implication in
Lemma~2.1, including the tail estimate. No compactness of the unknown zero
set beyond its proved discreteness and inverse-zero summability was assumed.
The positivity claim itself was \emph{not} supplied. Positivity of the theta
measure proves a different moment inequality, and the toy example refutes
the proposed automatic transfer to logarithmic moments. Replacing
$p(\lambda)^2$ by $|p(\lambda)|^2$, identifying two reflected factors without
justification, or inferring all-degree positivity from the table would each
introduce an invalid essential step. In general, nonnegative leading
principal determinants alone are not a test for positive semidefiniteness;
the condition used here is nonnegativity of every quadratic form of each
finite matrix.

There is no dependency on GRH, the simplicity conjecture, or a conjectural
spectral operator. Nor does this argument establish them: the explicit
nonvanishing on the real strip and the zero correspondence used above are
specific to the chosen zeta problem. A complete proof along this route still
requires an arithmetic/theta identity or an inequality proving
\eqref{ra:2:missing-positivity} for every polynomial, with no degree-dependent
exceptions. A disproof would require a rigorously certified negative value
for this \emph{actual} moment functional, not for either model function.

\Needspace{11\baselineskip}
\paragraph{Outcome.}
\textbf{Outcome: Exploratory approach with unresolved gap.}

The notebook supplies an exact moment-positivity reduction with a constructive
negative-polynomial witness under a hypothetical off-line zero, a quantitative
finite-test perturbation obstruction, and explicit failures of two tempting
positivity/symmetry arguments. It also isolates a specific incorrect inference
in one of the catalogue's cited proof claims. These results are proved or
qualified locally; they are not presented as historically new theorems about
zeta zeros. The unresolved step is all-degree positivity of the actual
logarithmic-moment functional. No complete proof or counterexample candidate
for the Riemann Hypothesis has been obtained in this batch.
% END RESEARCH ATTEMPT -- batch 1, problem 2


\subsection{Sources}
\begin{itemize}\raggedright
\item[\textnormal{[S1]}]\hypertarget{source:2:1}{} Enrico Bombieri, Problems of the Millennium: the Riemann Hypothesis, Section I and the prime-counting equivalence.
\url{https://www.claymath.org/wp-content/uploads/2022/05/riemann.pdf}.
{\small\textit{Catalogue formulation source. Consulted 24 September 2026: Section I, definition and zero-location statement.}}
\item[\textnormal{[S2]}]\hypertarget{source:2:2}{} Riemann Hypothesis.
\url{https://www.claymath.org/millennium/riemann-hypothesis/}.
{\small\textit{Catalogue research/status reference; not a proof endorsement.}}
\item[\textnormal{[S3]}]\hypertarget{source:2:3}{} A new proof that more than 2/3 of the zeros of the Riemann zeta function are simple and on the critical line.
\url{https://arxiv.org/abs/2609.02882}.
{\small\textit{Catalogue research/status reference; not a proof endorsement.}}
\item[\textnormal{[S4]}]\hypertarget{source:2:4}{} Lamzouri lecture at Max Planck Institute for Mathematics, September 3, 2026.
\url{https://math-events.uni-bonn.de/event/1560/}.
{\small\textit{Catalogue research/status reference; not a proof endorsement.}}
\item[\textnormal{[S5]}]\hypertarget{source:2:5}{} On the Extended, Generalized, and Grand Riemann Hypotheses: A Unified Approach Based on the General Properties of L-Functions — Weicun Zhang, v19.
\url{https://www.preprints.org/manuscript/202506.0481}.
{\small\textit{Catalogue research/status reference; not a proof endorsement.}}

% BEGIN ADDED RESEARCH SOURCES -- batch 1, problem 2
\item[\textnormal{[E1]}]\hypertarget{extra:2:1}{} Dave Platt and Tim Trudgian,
\emph{The Riemann hypothesis is true up to $3\cdot10^{12}$},
arXiv:2004.09765v1 (21 April 2020), Theorem~1 and Section~2;
published in \emph{Bulletin of the London Mathematical Society} (2021).
\url{https://arxiv.org/pdf/2004.09765v1}.
{\small\textit{Research reference: rigorous bounded-height verification,
including completeness of the zero count. Consulted 24 September 2026.}}

\item[\textnormal{[E2]}]\hypertarget{extra:2:2}{} Larry Guth and James Maynard,
\emph{New large value estimates for Dirichlet polynomials},
\emph{Annals of Mathematics} 203 (2026), no.~2, 623--675,
DOI: 10.4007/annals.2026.203.2.6, Theorem~1.2.
\url{https://annals.math.princeton.edu/2026/203-2/p06}.
Author manuscript: arXiv:2405.20552v2 (7 April 2026),
\url{https://arxiv.org/pdf/2405.20552v2}.
{\small\textit{Research reference: zero-density estimate; publication metadata
and the stated theorem were checked. Consulted 24 September 2026.}}

\item[\textnormal{[E3]}]\hypertarget{extra:2:3}{} Michael Griffin, Ken Ono,
Larry Rolen and Don Zagier,
\emph{Jensen polynomials for the Riemann zeta function and other sequences},
arXiv:1902.07321v2 (31 March 2019), Section~1, Theorems~1--2;
published in \emph{Proceedings of the National Academy of Sciences} (2019).
\url{https://arxiv.org/pdf/1902.07321v2}.
{\small\textit{Research reference: the fixed-degree, sufficiently-large-shift
quantifiers, not all-degree hyperbolicity. Consulted 24 September 2026.}}

\item[\textnormal{[E4]}]\hypertarget{extra:2:4}{} Brad Rodgers and Terence Tao,
\emph{The de Bruijn--Newman constant is non-negative},
\emph{Forum of Mathematics, Pi} 8 (2020), e6,
DOI: 10.1017/fmp.2020.6.
Author manuscript arXiv:1801.05914v5 (3 July 2021),
Section~1, equations~(1)--(4), and the main theorem.
\url{https://arxiv.org/pdf/1801.05914v5}.
{\small\textit{Research reference: the theta-kernel normalization and the
endpoint obstruction. Consulted 24 September 2026.}}

\item[\textnormal{[E5]}]\hypertarget{extra:2:5}{} David W. Farmer,
\emph{Jensen polynomials are not a plausible route to proving the Riemann
Hypothesis}, arXiv:2008.07206, Sections~2--4.
\url{https://arxiv.org/pdf/2008.07206}.
{\small\textit{Research reference: differentiation and information loss;
the paper's methodological assessment is not used as an impossibility theorem.
Consulted 24 September 2026.}}

\item[\textnormal{[E6]}]\hypertarget{extra:2:6}{} Jonathan Holland,
\emph{A new hyperbolicity wedge and a joint semicircle limit for Jensen
polynomials of Riemann's $\xi$-function},
arXiv:2608.08682v1 (9 August 2026), Theorem~1.1.
\url{https://arxiv.org/html/2608.08682v1}.
{\small\textit{Research reference: recent manuscript, statement inspected,
full proof not independently audited here; not used as a premise of any lemma.
Consulted 24 September 2026.}}

\item[\textnormal{[E7]}]\hypertarget{extra:2:7}{} Kiran S. Kedlaya,
\emph{More on the zeroes of $\zeta$}, MIT 18.785: Analytic Number Theory,
spring 2007, Section~1, Lemma~1 and Theorems~2--3.
\url{https://kskedlaya.org/18.785/zeroes.pdf}.
{\small\textit{Research reference: growth of the completed function, zero
counting, and Hadamard factorization. Consulted 24 September 2026.}}

\item[\textnormal{[E8]}]\hypertarget{extra:2:8}{} Paul Garrett,
\emph{12. Weierstrass and Hadamard products}, author lecture notes,
3 February 2021, Section~3 (Hadamard products).
\url{https://www-users.cse.umn.edu/~garrett/m/complex/notes_2020-21/12_Hadamard_products.pdf}.
{\small\textit{Research reference: canonical products and the finite-order
factorization theorem. Consulted 24 September 2026.}}

\item[\textnormal{[E9]}]\hypertarget{extra:2:9}{} Youness Lamzouri,
\emph{A new proof that more than $2/3$ of the zeros of the Riemann zeta
function are simple and on the critical line},
arXiv:2609.02882v2 (8 September 2026), Theorem~1.1,
Proposition~2.1 and Sections~2--3.
\url{https://arxiv.org/html/2609.02882v2}.
{\small\textit{Version-specific research supplement to [S3]: the stated bound
and method were inspected. Neither a complete analytic proof audit nor a
formal-certificate audit was performed. Consulted 24 September 2026.}}
% END ADDED RESEARCH SOURCES -- batch 1, problem 2
\end{itemize}

\end{document}
